% Capítulo 1 - Introdução  (PT-PT; espelho de ../thesis/ch1/chapter1.tex)
% Estrutura canónica MEIA. Tradução pura: labels/refs/citações idênticos; só a língua muda.

\chapter{Introdução}
\label{chap:introduction}

\section{Contextualização}
\label{sec:intro_context}

A maioria dos norte-americanos possui hoje ações, diretamente ou através de contas de reforma e fundos de
investimento \autocite{gallup2025stock}, e investe no maior mercado acionista do mundo, avaliado em
62,2~biliões de dólares no final de 2024 \autocite{sifma2025factbook} (Figura~\ref{fig:us_market_cap}).
A negociação decorre sobretudo na \gls{NYSE} e na \gls{NASDAQ}, e uma fração crescente vem agora de pessoas
comuns que usam aplicações móveis sem comissões, e não de instituições profissionais.

\begin{figure}[ht]
\centering
\includegraphics[width=0.9\textwidth]{figures/us_equity_market_cap.pdf}
\caption[Capitalização do mercado acionista dos EUA, 2015--2024]{Capitalização do mercado acionista dos
EUA, 2015--2024 (biliões de dólares). Elaborada a partir dos dados reportados no SIFMA 2025 Capital
Markets Fact Book \autocite{sifma2025factbook}.}
\label{fig:us_market_cap}
\end{figure}

Estes investidores não profissionais, ou ``de retalho'', são o público deste trabalho, e a sua posição é
difícil. Os mercados movem-se continuamente ao longo de milhares de títulos, e chegam notícias relevantes
durante todo o dia de negociação, mas os investidores de retalho não têm nenhuma das ferramentas que os
bancos e os fundos tomam como garantidas. A posse também é desigual: chega a 87\% dos agregados que ganham
acima de 100\,000 dólares, mas apenas a 28\% dos que ganham abaixo de 50\,000 dólares
\autocite{gallup2025stock}.

Entretanto, a inteligência artificial espalhou-se depressa pelas finanças: um inquérito de 2026 encontrou
81\% das empresas financeiras a adotar \gls{AI} e 71\% já a usar IA generativa \autocite{ccaf2026aifs}. Mas
grande parte desta tecnologia é opaca, raramente mostrando como chega a uma conclusão
\autocite{arrieta2020xai, adadi2018peeking}. Para quem tem de agir sobre um alerta e viver com o resultado,
essa opacidade é o obstáculo que esta dissertação remove: permite ao investidor ver \emph{porquê} um alerta
foi emitido.

\section{Definição do Problema}
\label{sec:intro_problem}

Quando uma posição se mexe, o investidor de retalho faz sempre as mesmas três perguntas, pela mesma
ordem, e elas correspondem a três problemas técnicos distintos:

\begin{enumerate}
  \item \textbf{Isto é realmente invulgar, ou é normal para esta ação?} Um dia de 3\% é banal numa
        empresa e um acontecimento raro noutra, pelo que o juízo é estatístico e tem de ser feito
        contra o comportamento recente de cada título.
  \item \textbf{É esta empresa, ou é o mercado todo?} Uma queda que se limita a acompanhar um índice
        em queda carrega informação diferente de uma queda que a empresa produziu por si, e separar
        as duas exige estimar com que força o título costuma seguir o mercado e o seu setor.
  \item \textbf{Já aconteceu algo assim, e o que se seguiu?} Responder a isto implica encontrar
        eventos passados genuinamente comparáveis, o que é um problema de recuperação sobre
        significado e não sobre palavras-chave, e depois medir o que os preços fizeram a seguir.
\end{enumerate}

As ferramentas gratuitas de consumo não respondem a nenhuma das três. Reportam a variação
percentual, que é o ponto de partida da primeira pergunta e não a sua resposta. Os terminais
profissionais respondem às três, a um custo de subscrição muito acima do alcance de um investidor
não-profissional. A lacuna que esta dissertação aborda não é, portanto, falta de dados, que são
abundantes e em larga medida gratuitos, mas falta de \emph{contexto} que um investidor comum
consiga obter e, sobretudo, consiga verificar.

\subsection*{A quem se destina}

Dois perfis de utilizador enquadram o desenho, e querem coisas opostas do mesmo sistema.

O primeiro é o \textbf{detentor de longo prazo}: alguém com um número reduzido de posições, que as
consulta com pouca frequência e sente ansiedade quando o ecrã está vermelho. Para esta pessoa, a
saída mais valiosa é muitas vezes a permissão para não fazer nada, entregue como evidência de que o
movimento foi generalizado ao mercado e nada teve de notável. O modo de falha a evitar é o alarme
sem informação.

O segundo é o \textbf{investidor de retalho ativo}: alguém que consulta várias vezes por dia e já
recebe alertas que apenas dizem que um preço se mexeu. Para esta pessoa, o valor está no contexto
chegar com o alerta em vez de depois de uma pesquisa noutro lado. O modo de falha a evitar é a
fadiga de alertas, em que o volume ensina o leitor a ignorar o canal.

Os dois perfis puxam em direções diferentes quanto à cadência, e o sistema resolve essa tensão
sendo deliberadamente seletivo e enunciando as suas regras de seleção de forma aberta, decisão
retomada no Capítulo~\ref{chap:investigator}.

\subsection*{Posição assumida}

Este trabalho assume uma postura deliberadamente diferente da das ferramentas que o rodeiam: não
prevê preços, \emph{explica} eventos. Para cada movimento abrupto ou notícia relevante, o sistema
mostra o que detetou, como, de onde veio a informação, e quais os eventos passados que se assemelham
ao atual, juntamente com o que aconteceu aos preços depois (Figura~\ref{fig:concept}).

Essa recusa é uma restrição de desenho e não uma limitação pela qual pedir desculpa. Exclui importar
previsões de terceiros, o que contradiria a afirmação que está no centro do sistema, e torna cada
saída verificável, porque uma afirmação sobre o que já aconteceu pode ser confrontada com o registo,
ao passo que uma previsão não pode.

\begin{figure}[ht]
\centering
\begin{tikzpicture}[
  font=\small, >={Stealth[length=2.2mm]},
  trig/.style={draw, rounded corners, fill=black!8, align=center, text width=2.7cm, minimum height=0.9cm},
  sys/.style={draw, rounded corners, fill=black!3, align=center, text width=2.2cm, minimum height=1.4cm},
  outbox/.style={draw, rounded corners, fill=black!3, align=center, text width=3.4cm, minimum height=1.4cm},
]
  \node[trig] (t1) {Movimento abrupto de mercado};
  \node[trig, below=6mm of t1] (t2) {Notícia relevante};
  \node[sys] (sys) at ($(t1)!0.5!(t2) + (3.6,0)$) {\textbf{InvestiGator}};
  \node[outbox, right=1.2cm of sys] (out) {Alerta explicado\\ \footnotesize evento $\cdot$ raciocínio\\ \footnotesize fontes $\cdot$ precedentes};
  \node[right=0.9cm of out] (inv) {\footnotesize investidor};
  \draw[->] (t1) -- (sys);
  \draw[->] (t2) -- (sys);
  \draw[->] (sys) -- (out);
  \draw[->] (out) -- (inv);
\end{tikzpicture}
\caption[O que o InvestiGator faz]{Os dois gatilhos, um movimento súbito de mercado ou uma notícia
relevante, alimentam o InvestiGator, que responde com um alerta explicado que o investidor pode seguir e
verificar, e não apenas uma notificação seca.}
\label{fig:concept}
\end{figure}

\section{Questões de Investigação e Objetivos}
\label{sec:intro_objectives}

O objetivo é conceber, construir e avaliar um sistema de alertas financeiros explicável para investidores de
retalho dos EUA, movido por dois gatilhos independentes, um movimento súbito de mercado e uma notícia
recebida, e entregue por Telegram. Quatro questões de investigação guiam o trabalho:

\begin{itemize}
  \item \textbf{\gls{RQ}1.} Conseguem métodos estatísticos transparentes detetar movimentos abruptos de
        mercado de forma suficientemente consistente para serem úteis a um investidor de retalho,
        mantendo-se explicáveis?
  \item \textbf{\gls{RQ}2.} Consegue um motor de correlação notícia--mercado recuperar notícias históricas
        genuinamente análogas melhor do que baselines triviais, e quantificar o impacto observado nos preços
        desses precedentes sem lookahead?
  \item \textbf{\gls{RQ}3.} Podem as explicações resultantes ser fiéis ao que o sistema realmente calcula, e
        úteis a um investidor de retalho?
  \item \textbf{\gls{RQ}4.} Consegue um modelo supervisionado, treinado em notícias históricas e contexto de
        mercado, priorizar utilmente que notícias merecem um alerta, para além de baselines simples de
        volatilidade, sem nunca prever a direção do preço?
\end{itemize}

Estas questões traduzem-se em cinco objetivos concretos:
\begin{itemize}
  \item construir um pipeline de alertas explicável, movido pelos dois gatilhos;
  \item avaliar o detetor de anomalias face a baselines transparentes;
  \item avaliar a recuperação de notícias históricas face a baselines triviais;
  \item aferir se as explicações são fiéis e úteis;
  \item treinar, calibrar e avaliar honestamente um modelo de triagem de materialidade que ordena os alertas
        de notícias.
\end{itemize}
O sistema é construído de forma incremental, primeiro na sua forma mais simples e defensável, uma escolha
metodológica explicada no Capítulo~\ref{chap:methods}.

\section{Contribuições Científicas}
\label{sec:intro_contributions}

Esta é uma dissertação em \emph{Engenharia} de Inteligência Artificial: a contribuição não é um algoritmo
novo, mas a montagem disciplinada de algoritmos existentes num sistema que funciona, se explica a si próprio
e pode ser reproduzido. Quatro contribuições destacam-se:

\begin{itemize}
  \item Um método reprodutível para relacionar notícias com o impacto no mercado, que junta recuperação por
        embeddings de frases com janelas de estudo de evento, com a medida de similaridade e os horizontes
        temporais declarados de forma explícita e sem deixar entrar qualquer informação sobre o futuro.
  \item O \textbf{InvestiGator}, um sistema de alertas que mostra toda a sua cadeia de raciocínio: o evento
        detetado, a evidência de apoio, as fontes e os precedentes históricos relevantes. Como cada
        explicação é construída diretamente a partir das quantidades que o sistema calcula, o texto não pode
        divergir do cálculo.
  \item Um modelo de \emph{triagem de materialidade} treinado e calibrado: aprendizagem supervisionada sobre
        notícias históricas e contexto de mercado que estima com que frequência notícias com um dado perfil
        foram seguidas de um movimento anormal, rotuladas pelo próprio código de estudo de evento do sistema,
        avaliadas face a uma baseline de volatilidade, e integradas nos alertas como evidência ordenada e
        explicada, e não como uma previsão.
  \item Uma avaliação crítica e honesta de cada componente central em dados reais, face a baselines simples e
        lexicais, com os resultados, as ablações e as limitações todos reportados com clareza.
\end{itemize}

\section{Estrutura do Documento}
\label{sec:intro_structure}

O resto da dissertação foi escrito para ser lido por ordem, cada capítulo respondendo a uma pergunta:

\begin{itemize}
  \item \textbf{Capítulo~\ref{chap:sota} (Estado da Arte):} o que já se sabe, e o que escolhe este trabalho?
  \item \textbf{Capítulo~\ref{chap:methods} (Métodos e Materiais):} como é construído o sistema, e como é
        avaliado?
  \item \textbf{Capítulo~\ref{chap:investigator} (InvestiGator):} o que é o sistema, e como se encaixam os
        seus dados, partes, fluxo e decisões?
  \item \textbf{Capítulo~\ref{chap:casestudies} (Estudos de Caso):} funciona em dados reais?
  \item \textbf{Capítulo~\ref{chap:conclusions} (Conclusões):} o que se aprendeu, e o que vem a seguir?
\end{itemize}
